\documentclass[UTF8,a4paper,11pt]{ctexart}

\usepackage{amsmath,amssymb,bm,mathtools}
\usepackage{booktabs}
\usepackage{geometry}
\usepackage{graphicx}
\usepackage{xcolor}
\usepackage{enumitem}
\usepackage{cite}
\usepackage{hyperref}

\geometry{left=24mm,right=24mm,top=25mm,bottom=25mm}
\setlength{\parindent}{2em}
\setlength{\parskip}{0.25em}
\hypersetup{
  colorlinks=true,
  linkcolor=blue!55!black,
  citecolor=green!40!black,
  urlcolor=blue!60!black
}

\newcommand{\R}{\mathbb{R}}
\newcommand{\Prob}{\mathbb{P}}
\newcommand{\norm}[1]{\left\lVert #1 \right\rVert}
\newcommand{\pos}[1]{\left[#1\right]_{+}}
\newcommand{\argmin}{\operatorname*{arg\,min}}

\title{面向遮挡物操作的阶段条件语义预测与图像空间视觉伺服\\
\large 从自动语义标注到可验证闭环控制的中文方法草稿}
\author{作者待定}
\date{\today}

\begin{document}

\maketitle

\begin{abstract}
动作分块 Transformer（ACT）能够从示范中学习复杂操作，但长动作块在执行期间可能因视觉误差、接触扰动和执行器滞后而逐渐过时。本文提出一个精简的闭环框架：首先利用人工审核过的稀疏提示、SAM2 和轻量级 YOLO 提示生成器自动构建语义标注，并对自动标签进行离线质量评估；随后从任意数量的相机视角提取面积、质心、距离和接触等可解释语义状态；利用带顺序约束的时序模型自动学习任务阶段；以 ACT 生成名义动作块，以动作条件语义动力学预测动作结果，并通过图像空间视觉伺服对名义动作进行小幅闭环修正。该方法不强制要求相机标定：单视角或多视角均可在各自归一化图像坐标中工作；只有需要将多个视角融合到统一物理坐标、报告实际距离或声明物理空间控制保证时才需要相机标定。本文给出方法成立所需的可检查条件，并将结论限定为已验证运行域内的局部可靠性，而非任意初态和任意遮挡下的全局成功保证。
\end{abstract}

\noindent\textbf{关键词：}自动语义标注；动作分块 Transformer；阶段学习；语义预测；无标定视觉伺服；不确定性评估

\section{问题与核心思路}

任务由机器人状态 $\bm q_t\in\R^{n_q}$、可用相机集合 $\mathcal V_t$ 和动作 $\bm a_t$ 描述。相机数量不固定：$|\mathcal V_t|=1$ 时为单视角，$|\mathcal V_t|>1$ 时为多视角；某个相机暂时失效时，只要仍有可靠视角，系统可以继续工作。

本文保留四个核心模块：
\begin{enumerate}[leftmargin=2.3em]
  \item 自动生成语义标注，并在训练前评估每帧、每类标签质量；
  \item 从一个或多个视角构造可解释语义状态，并自动学习任务阶段；
  \item ACT 预测名义动作，语义动力学预测这些动作将产生的视觉结果；
  \item 图像空间视觉伺服根据实际语义反馈修正 ACT，并在观测不可靠时缩短执行或重新规划。
\end{enumerate}

该框架的目的不是用手工控制器替代 ACT。ACT 负责复杂的双臂协同和接触先验，语义闭环只负责发现动作块已经过时并做受限修正。

\section{自动语义标注与标签不确定性}

\subsection{从少量人工审核到完整数据集}

语义类别包括遮挡物或布料、目标物体、目标区域和工具/执行器，分别记为 $C,O,G,A$。完整标注按照以下流程生成：
\begin{enumerate}[leftmargin=2.3em]
  \item 使用 Codex 为少量代表视频生成初始离散 JSON。每条记录包含视频、帧号、类别和包围盒，可直接作为分割提示；
  \item 将 JSON 渲染为带包围盒和类别名称的图片，由人工审核并修正错误框。人工审核的是少量关键提示，而不是逐帧描绘 mask；
  \item 将审核后的框输入 SAM2，在少量种子视频中传播得到完整时序语义 mask\cite{ravi2024sam2}；
  \item 使用这些完整种子标注训练轻量级 YOLO 网络\cite{redmon2016yolo}，使其为其余视频自动生成类别框提示；
  \item 将 YOLO 提示再次交给 SAM2 传播，生成整个训练数据集的逐帧语义标注。
\end{enumerate}

这个流程显著减少人工成本，但最终标签仍是自动生成的。Codex 初始框、YOLO 框和 SAM2 时序传播中的任一环节都可能产生系统误差，因此标签在进入策略训练前必须经过质量门。

\subsection{离线标签质量评估}

对视角 $v$、类别 $c$ 和二值标签 $M_{t,c}^{v}$，记录以下原始诊断量：

\paragraph{时间一致性}
\begin{equation}
q_{iou,t}^{v,c}=\frac{1}{2}\left[
\operatorname{IoU}(M_{t-1,c}^{v},M_{t,c}^{v})+
\operatorname{IoU}(M_{t,c}^{v},M_{t+1,c}^{v})\right].
\end{equation}

\paragraph{面积突变}
令 $r_{t,c}^{v}=|M_{t,c}^{v}|/|\Omega_v|$，定义
\begin{equation}
d_{area,t}^{v,c}=\left|
\log\frac{r_{t,c}^{v}+\epsilon}{r_{t-1,c}^{v}+\epsilon}
\right|.
\end{equation}

\paragraph{质心跳变}
将质心坐标除以图像宽高，得到归一化质心 $\bm p_{t,c}^{v}\in[0,1]^2$，定义
\begin{equation}
d_{center,t}^{v,c}=\frac{
\norm{\bm p_{t,c}^{v}-\bm p_{t-1,c}^{v}}
}{\sqrt{2}}.
\end{equation}

\paragraph{连通域异常}
统计去除极小噪声区域后的有效连通域数量 $n_{cc,t}^{v,c}$。布料可能被机械臂遮挡而自然分裂，工具也可能包含多个部件，因此面积、运动和连通域阈值必须按“视角--类别”设置，不能所有类别、所有相机共用一个阈值。例如本任务中侧视图的工具容易与背景遮挡布混淆，而正视图受此影响较小，二者必须分别校准。

将各指标归一化为 $[0,1]$ 的子分数后，计算
\begin{equation}
Q_{t,c}^{v}=\frac{\sum_k w_k q_{t,c,k}^{v}}
{\sum_k w_k},\qquad Q_{t,c}^{v}\in[0,1].
\end{equation}
单次低 IoU 可能来自真实快速运动，不能直接等同于标签错误。系统保留所有原始指标供人工抽查，并使用在独立审核集上校准的综合分数产生有效标记
\begin{equation}
\gamma_{t,c}^{v}=\mathbb I[Q_{t,c}^{v}\geq\tau_c].
\end{equation}
训练时，低质量类别不参与当前分割损失；未来任一必要标签不可靠时，相应时间点不参与语义动力学损失。

离线标签质量与部署时模型置信度是两个不同问题。前者回答“训练标签是否可信”，后者回答“策略当前看到的分割是否可信”，二者需要分别评估。

\section{可变视角语义状态}

\subsection{每个视角的图像空间状态}

分割模型输出软概率 $P_{t,c}^{v}(x,y)\in[0,1]$。每个视角独立计算：
\begin{itemize}[leftmargin=2.3em]
  \item 各类面积比例 $r_{t,c}^{v}$；
  \item 物体、目标区域和执行器的归一化软质心；
  \item 物体到目标、执行器到物体的归一化图像距离；
  \item 物体--目标和物体--布料的图像接触指标；
  \item 分割熵、时间跳变和类别可见性等在线置信度。
\end{itemize}

当前语义图采用互斥多类 softmax，所以物体和目标区域的硬标签不会在同一像素直接重叠。第一版使用局部膨胀后的接触代理
\begin{equation}
\eta_{OG,t}^{v}=
\frac{\sum_{x,y}P_{t,O}^{v}(x,y)
\mathcal D_{\rho}(P_{t,G}^{v})(x,y)}
{\sum_{x,y}P_{t,O}^{v}(x,y)+\epsilon},
\end{equation}
其中 $\mathcal D_{\rho}$ 是局部软膨胀。它表示图像邻接或占据程度，不等于真实三维物理重叠。

将上述量组成视角状态 $\bm z_t^v$。全部量都在各自图像坐标内归一化，因此不需要知道相机内参、外参或桌面单应矩阵。

\subsection{单视角与多视角统一处理}

对每个可靠视角使用共享编码器 $\phi$，再进行与视角数量无关的集合聚合：
\begin{equation}
\bm h_t=\rho\left(
\frac{\sum_{v\in\mathcal V_t}\alpha_t^v
\phi(\bm g_t^v\odot\bm z_t^v,\bm g_t^v,e_v)}
{\sum_{v\in\mathcal V_t}\alpha_t^v+\epsilon}
\right),
\end{equation}
其中 $\bm g_t^v\in[0,1]^{d_z}$ 是按语义类别和特征计算的置信门，$\alpha_t^v$ 是该视角对当前任务相关特征的总体可用度，$e_v$ 是可学习的视角标识。这样可以只降低侧视图 tool 特征的权重，同时保留侧视图 object 和 region 的信息，而不是丢弃整个侧视角。一个视角时公式自然退化为单视角；多个视角时，不同相机提供遮挡互补信息。训练时随机丢弃视角和单个类别特征，使模型学会在观测缺失时退化运行。

未经标定时，不直接平均不同相机中的质心坐标，而是先分别编码再聚合。若任务需要厘米级距离、统一桌面轨迹或把多个视角严格融合到同一物理坐标，则需要相机内外参或桌面单应标定。这是物理坐标融合的条件，不是图像空间视觉伺服的必要条件。

\section{自动阶段学习}

\subsection{阶段定义}

任务阶段记为
\begin{equation}
\mathcal M=\{U,E,T,R,D\},
\end{equation}
分别表示掀开遮挡物、暴露并分离目标、运输目标、恢复遮挡物和完成。

单帧语义标签不能唯一决定阶段。例如开始和完成时都可能看到较大布料面积。因此阶段模型必须使用最近 $L$ 帧的语义状态、机器人状态和动作历史：
\begin{equation}
\bm\pi_t^m=p_\psi\left(
m_t\mid \bm h_{t-L+1:t},\bm q_{t-L+1:t},
\bm a_{t-L:t-1}\right).
\end{equation}

\subsection{不逐条人工标边界的学习方法}

五阶段可以主要通过学习获得，不必人工校正每条轨迹。训练流程为：
\begin{enumerate}[leftmargin=2.3em]
  \item 根据高质量语义序列自动产生软事件候选：物体持续可见、物体与布料持续分离、物体持续进入目标区域、布料恢复且动作趋于停止；
  \item 将事件候选转换为带置信度的软阶段标签，而不是立即生成不可修改的硬标签；
  \item 使用只允许 $U\rightarrow E\rightarrow T\rightarrow R\rightarrow D$ 前向转移的 HMM、GRU 或 Transformer 学习阶段；
  \item 对低置信、阶段逆序、事件缺失和模型分歧大的轨迹进行人工抽查；
  \item 用审核过的小子集校准阈值和阶段概率，再重新生成全部阶段标签。
\end{enumerate}

完全无人工锚点也可以聚类出若干时序状态，但不能保证“第 2 类”在语义上一定对应运输阶段，这属于阶段置换和可辨识性问题。因此合理目标是“自动标注全部轨迹，只审核少量困难轨迹”，而不是“完全取消任何人工验证”。

部署时对阶段概率使用连续帧投票、滞回和最小驻留时间，防止一次分割抖动造成阶段来回切换。

\section{语义预测 ACT 与无标定视觉伺服}

\subsection{ACT 名义动作}

ACT 根据 RGB、软语义图、机器人状态和阶段表示预测长度为 $H$ 的名义动作块\cite{zhao2023act}：
\begin{equation}
\bm a_{t:t+H-1}^{ACT}=\pi_{ACT}(I_t,\bm h_t,\bm q_t,\bm\pi_t^m).
\end{equation}
ACT 保留复杂示范行为，策略以较短间隔重新观察和规划，而不是盲目执行完整长动作块。

\subsection{动作条件语义动力学}

学习模型预测给定动作后的未来语义状态及不确定性：
\begin{equation}
p_\theta\left(
\bm h_{t+k}\mid \bm h_t,\bm q_t,
\bm a_{t:t+k-1},\bm\pi_t^m
\right),\qquad k\in\mathcal K.
\end{equation}
预测均值用于判断视觉进度，预测方差和独立校准残差共同构成可信区间。训练 softmax 自信并不代表动力学误差有界，因此误差界必须在未参与训练的轨迹上校准。

\subsection{图像空间视觉目标}

每个阶段定义简单、可解释的图像目标：
\begin{align}
V_U &= \pos{r_C-\bar r_U}^2+\lambda_{Uv}\pos{\bar r_O-r_O}^2,\\
V_E &= \lambda_{Ec}\pos{\eta_{OC}-\bar\eta_{OC}}^2,\\
V_T &= \pos{d_{OG}-\bar d_T}^2+
       \lambda_{T\eta}\pos{\bar\eta_T-\eta_{OG}}^2,\\
V_R &= \lambda_{Rc}\pos{\bar r_R-r_C}^2+
       \lambda_{R\eta}\pos{\bar\eta_D-\eta_{OG}}^2,\\
V_D &= 0.
\end{align}
单视角时直接使用该视角的 $V_m^v$；多视角时按可靠度加权，或对关键误差采用保守最大值。阈值由成功训练轨迹和独立校准轨迹确定，不能在最终测试集上反复调整。

\subsection{为什么不一定需要相机标定}

经典基于图像的视觉伺服可以直接让图像误差下降\cite{chaumette2006visual}。本文不把像素速度直接解释为三维末端速度，而是从示范数据学习局部关系
\begin{equation}
\Delta\bm h_t=\bm f_m(\bm h_t)+
\bm G_m(\bm h_t)\Delta\bm a_t+\bm d_t,
\end{equation}
其中 $\Delta\bm a_t$ 是关节目标增量。只要局部模型 $\bm G_m$ 在当前运行域内准确，便可以在图像空间判断某个关节修正是否使 $V_m$ 下降，不需要显式相机标定。这属于学习型、无标定图像视觉伺服。

相机标定在以下情况才是必要或明显有益的：
\begin{itemize}[leftmargin=2.3em]
  \item 要把不同视角的点融合成统一三维点或桌面坐标；
  \item 要直接使用解析图像雅可比矩阵，而不是学习局部动作--图像关系；
  \item 要报告厘米、毫米等物理尺度误差；
  \item 要把碰撞距离、工作空间边界等物理安全约束写入控制器。
\end{itemize}

\subsection{受限闭环修正}

令 ACT 当前动作增量为 $\Delta\bm a_t^{ACT}$。在局部模型上寻找距离 ACT 最近的修正：
\begin{align}
\Delta\bm a_t^*,\delta_t^*=\argmin_{\Delta\bm a,\delta\geq0}\quad
&\frac{1}{2}\norm{\Delta\bm a-\Delta\bm a_t^{ACT}}_{W}^{2}
+p_\delta\delta^2\\
\text{s.t.}\quad
&\widehat{\Delta V_m}(\bm h_t,\Delta\bm a)
+\beta_{model}+\beta_{obs}
\leq-\rho_m V_m(\bm h_t)+\delta,\\
&\Delta\bm a_{min}\leq\Delta\bm a\leq\Delta\bm a_{max}.
\end{align}
这里 $\beta_{model}$ 和 $\beta_{obs}$ 分别上界动力学预测误差和观测误差。QP 只做满足视觉下降条件所需的最小修正；若松弛 $\delta$ 过大或观测不可靠，则不继续声称下降，而是停止、减速或立即重新规划。

\section{训练目标与执行流程}

总损失保留四项：
\begin{equation}
\mathcal L=\mathcal L_{ACT}
+\lambda_{seg}\mathcal L_{seg}^{quality}
+\lambda_{phase}\mathcal L_{phase}
+\lambda_{dyn}\mathcal L_{dyn}^{quality}.
\end{equation}
其中分割和动力学损失仅在质量标记 $\gamma=1$ 的标签上计算。阶段损失使用软伪标签与少量审核标签；动力学损失预测未来语义分布，而非生成完整 RGB。

推荐按以下顺序训练：
\begin{enumerate}[leftmargin=2.3em]
  \item 完成自动语义标注、离线质量评估和小规模人工校准；
  \item 训练分割与 ACT 基线，并验证单视角和多视角输入；
  \item 自动生成阶段软标签，训练时序阶段模型并审核困难轨迹；
  \item 使用专家动作训练语义动力学，在独立校准集上估计误差界；
  \item 先以影子模式记录 QP 修正，不发送给机器人；验证后再逐步开放小幅修正。
\end{enumerate}

在线控制循环为：获取可用视角、预测语义和置信度、更新阶段、ACT 预测动作块、语义模型预测未来、QP 计算首个受限修正、执行少量动作，然后重新观察。若视角置信度过低、阶段异常变化、实际语义超出预测区间或 $V_m$ 未按预期下降，则提前重新规划。

\section{方法成立的条件}

\subsection{局部可靠性结论}

在阶段 $m$ 的已验证运行域 $\mathcal D_m$ 内，假设真实离散视觉变化与学习模型的误差对 $V_m$ 的影响不超过 $\beta_{model}+\beta_{obs}$。若 QP 可行并满足上述鲁棒下降约束，则实际闭环满足
\begin{equation}
V_m(\bm h_{t+1})-V_m(\bm h_t)
\leq-\rho_mV_m(\bm h_t)+\delta_t.
\end{equation}
当 $\delta_t\leq\bar\delta_m$ 时，$V_m$ 收敛到大小与 $\bar\delta_m/\rho_m$ 有关的邻域。该结论是局部、条件性的：它只在观测可信、模型误差界有效且 QP 可行时成立。

\subsection{条件是否能够达到}

\begin{table}[htbp]
  \centering
  \small
  \begin{tabular}{p{0.25\linewidth}p{0.43\linewidth}p{0.20\linewidth}}
    \toprule
    条件 & 如何检查或实现 & 可达到性 \\
    \midrule
    自动标签质量可控 & 人工审核独立小样本，校准各类质量分数，屏蔽低质量帧 & 可以，但不能零人工 \\
    至少一个视角可观测 & 训练视角丢弃，部署时使用置信门并允许降级运行 & 通常可以 \\
    阶段可辨识 & 使用时间历史、顺序约束、软标签和困难轨迹审核 & 可以近似达到 \\
    局部动作--视觉模型准确 & 划分独立校准集，检查预测区间覆盖率和域外检测 & 局部可以 \\
    视觉目标在动作限制内可控 & 离线检查 QP 可行率，实机影子测试，限制修正幅度 & 需实验验证 \\
    延迟和执行长度有界 & 短间隔重规划，记录端到端延迟，异常时中断 & 工程上可以 \\
    测试状态接近训练域 & 扩充初始位置、遮挡和恢复数据，并进行域外检测 & 只能部分满足 \\
    \bottomrule
  \end{tabular}
  \caption{可靠性条件、检查方法与现实可达到性。}
\end{table}

这些条件在受控工作台、固定任务对象和覆盖充分的数据中是可以近似达到并实证检查的，但很难在任意新物体、任意相机位置、完全遮挡或严重缠绕下始终成立。因此合理表述是“在校准运行域内提高闭环可靠性并满足局部视觉下降条件”，而不是“保证所有任务必然成功”。

\section{实验环境与验证方案}

建议首先在受控桌面操作环境中验证：机器人、任务对象和工作区域保持一致，使用一个或多个固定 RGB 相机，控制频率、图像延迟和动作执行间隔均可记录。相机可以有不同分辨率和观察方向，但训练阶段必须覆盖部署时可能出现的相机组合；若相机移动、镜头更换或视野明显改变，应视为新的运行域并重新校准质量阈值和动力学误差界。当前已知侧视图的 tool 会受到背景遮挡布干扰，因此实验必须单独报告 front-tool 与 side-tool 的标签质量、在线置信度和消融结果。

数据按完整轨迹划分为训练集、校准集和测试集，不能把同一轨迹的相邻帧分到不同集合。训练集用于分割、阶段模型、ACT 和语义动力学训练；校准集用于确定标签质量阈值、阶段置信阈值、预测区间和 QP 裕量；测试集只用于最终报告。为了验证单视角与多视角能力，应分别报告单个相机、全部相机、随机缺失相机和短时遮挡四种设置。

评估至少包含四组指标：自动标签在人工审核子集上的类别 IoU/Dice 与质量分校准误差；四个事件边界的平均帧误差、阶段准确率和低置信轨迹召回率；动作条件语义预测误差及预测区间覆盖率；实机任务成功率、重规划次数、QP 可行率和视觉误差实际下降比例。只有当这些指标在独立测试轨迹上稳定，且部署环境未明显超出训练与校准范围时，才认为前述局部可靠性条件近似成立。

\section{简单说明}

整个方法可以理解为：先用少量人工审核教会自动标注系统“画对语义 mask”，再给每张自动标签打质量分；策略只相信高质量标签。机器人运行时，ACT 先给出它根据示范学到的一段动作，语义预测模型判断这段动作大概会把物体和布料带到哪里。执行少量动作后，系统重新看图像：如果物体确实更接近目标，就继续；如果没有进展、阶段判断改变或视觉不可信，就缩短执行并重新规划。视觉伺服只在图像坐标中要求“误差变小”，所以单相机可以工作，多相机只是提供更多互补证据，并不天然要求相机标定。

\section{局限性}

\begin{enumerate}[leftmargin=2.3em]
  \item 自动质量指标只能发现时间跳变、面积异常、质心异常和拓扑异常，无法识别长期稳定但始终分错类别的系统偏差，因此仍需人工审核独立子集；
  \item 单视角完全遮挡时没有方法能从当前图像恢复真实状态，多视角只能降低而不能消除这个问题；当前侧视图 tool 与背景遮挡布的混淆属于已知的视角相关偏差，必须单独校准和降权；
  \item 无标定图像伺服只保证图像指标下降，不能直接解释为三维距离、碰撞距离或末端位姿误差下降；
  \item 接触代理来自互斥语义图的局部邻接，不是真实物理接触。更精确的判断需要 amodal 标注、深度或标定后的目标区域；
  \item 成功示范主要覆盖正常轨迹。卡住、推过目标、物体丢失和恢复行为需要额外失败/接管数据；
  \item 阶段伪标签可能继承语义错误。顺序约束能减少抖动，但不能修复系统性误标；
  \item 学习动力学和误差界只在校准域内有效。相机位置、光照、物体外观或机器人延迟显著变化后必须重新验证。
\end{enumerate}

\section{结论}

本文给出一个从自动语义数据到闭环策略的精简方案。方法以 Codex、人工审核、SAM2 和轻量 YOLO 构成低成本标注流水线，以离线质量门避免自动标签被无条件信任；以可变视角语义状态和时序模型自动学习任务阶段；以 ACT 保留复杂示范行为，以语义动力学和图像空间视觉伺服提供短周期反馈修正。相机标定不是图像空间闭环的必要条件，但在跨视角物理融合和物理安全约束中仍然必要。方法能够提供的是已验证运行域内、误差界成立时的局部可靠性，而不是不受条件限制的全局任务保证。

\bibliographystyle{unsrt}
\bibliography{references}

\end{document}
